\documentclass[12pt]{article}
\usepackage[margin=1in]{geometry}
\usepackage[T1]{fontenc}
\usepackage{xcolor}
\usepackage[hyphens]{url}
\usepackage{hyperref}
\hypersetup{hidelinks}
\definecolor{darkblue}{rgb}{0.1, 0.2, 0.6}

\newenvironment{comment}
{\begin{quotation}\noindent\small\itshape\color{darkblue}\ignorespaces}
{\end{quotation}}
\usepackage{booktabs,tabularx}
\usepackage[round]{natbib}   % enables \citep and \citet
\usepackage{placeins}        % \FloatBarrier keeps tables with their response

% Tables in this letter are numbered R1, R2, ... so they are not confused
% with the manuscript's own tables.
\renewcommand{\thetable}{R\arabic{table}}

% Let tall tables share a page with text instead of taking a page alone.
\renewcommand{\topfraction}{0.9}
\renewcommand{\bottomfraction}{0.8}
\renewcommand{\textfraction}{0.07}
\renewcommand{\floatpagefraction}{0.8}

\sloppy

\begin{document}

\begin{center}
{\Large\bfseries Response to the Comments}\\[0.5em]
{\large\itshape Principles for Open Data Curation:\\
A Case Study with the New York City 311 Service Request Data}
\end{center}

\bigskip

We thank the Editor and the Referee for their careful reading and
helpful comments, and we appreciate the Referee's positive assessment
of the paper. We have revised the manuscript in response; the main
changes are:

\begin{enumerate}
\item \textbf{Figures.} Figures have been revised, replaced, or
  removed to improve readability, with larger fonts, explicit units,
  and expanded captions that state each figure's key observations.
  One figure has been replaced by a table, and figures and tables have
  been renumbered accordingly.
\item \textbf{Durations.} We added an explanation of why many long
  durations are plausible given the agency and complaint type, and we
  now summarize negative durations numerically in the text.
\item \textbf{Data curation roles.} Following the Referee's
  suggestion, a new actionable step recommends formalizing a data
  curation responsibility within each agency.
\item \textbf{Clarity and corrections.} We added examples of redundant
  fields, removed a paragraph on embedded commas after verifying that
  the export handles them correctly, and corrected the typographical
  issues the Referee identified.
\end{enumerate}

Point-by-point responses follow, with the Referee's comments in
\emph{\color{darkblue}italic and blue}. Page and line numbers refer to the revised manuscript. Where
a figure has been renumbered, we give both its original and revised
numbers. Tables in this letter are numbered R1 and R2 to distinguish
them from tables in the manuscript.

% ===============================================================
\subsection*{To the Referee}
% ===============================================================

\FloatBarrier
\begin{comment}
Page 4, lines 31--32: The authors mention two potentially
redundant pairs. This is discussed later, but it might help to add
something like ``e.g., cross\_street and intersection\_street'' at this
point in the article.
\end{comment}

We thank the Referee for this suggestion, which improves the clarity of
the paragraph. We now give the \texttt{cross\_street} and
\texttt{intersection\_street} pairs as examples at first mention
(p.~4, ll.~10--12), which now reads: ``The dataset contains seven
street-related fields, including two likely redundant pairs, e.g.,
\texttt{cross\_street\_1} \& \texttt{intersection\_\allowbreak street\_1}
and \texttt{cross\_street\_2} \& \texttt{intersection\_\allowbreak street\_2}.''

\FloatBarrier
\begin{comment}
Page 4, bottom: This is just a comment for what it's worth.
I was doing text analysis on data like this that contained a Notes
field. It was a CSV file, and I imported it into MATLAB. I ended up
with a strange number of columns when importing it. It turns out that
the commas inside the Notes field were telling MATLAB that there were
multiple columns from the Notes. So, CSV can be dangerous or
problematic with free-form text fields. You are alluding to problems
like this at the bottom of the page. Just wanted to pass along my
experience.
\end{comment}

We thank the Referee for sharing this experience, which prompted us to
check the export directly. The NYC Open Data portal runs on the Socrata
platform, whose CSV exports quote any field containing commas, line
breaks, or quotation marks and escape embedded quotes by doubling
\citep{socrata_csv}, consistent with RFC~4180 \citep{rfc4180}. We scanned
the full, unmodified export (approximately 16.0 million records;
\citealp{nyc311_data}) and found embedded commas in 15 of the 41 fields
(see Table~\ref{tab:embedded-commas}, at the end of this letter). Apart from
\texttt{location}, which always contains one, over 2.6 million values
contain commas, mostly in the free-text \texttt{resolution\_description}
field, and about 163{,}000 contain escaped double quotes. All are correctly
quoted, and every record parses to the expected 41 columns with
\texttt{fread()} \citep{datatable2025}.

The Referee's point is well taken: a parser that ignores CSV quoting would
misread a substantial share of these records. Because the export is
correctly quoted and we read it with a quote-aware parser, the problem
does not arise with these data, and we have removed the paragraph on this
issue (formerly p.~4, ll.~45--46).

\FloatBarrier
\begin{comment}
Figure 2: Please add more explanation to the caption. What
do you want the reader to notice or take away from the figure? It's ok
to repeat what you have in the body of the paper.
\end{comment}

We thank the Referee for this suggestion. In revising the caption, we
concluded that the figure contained more detail than a reader could
readily absorb, and that its key points are conveyed more clearly in
the text. We have therefore removed the figure and revised the
accompanying discussion (p.~4, ll.~29--32) to state its main takeaway
directly: just six agencies account for 90.19\% of all SRs.

\FloatBarrier
\begin{comment}
Figure 3: First, this is very hard to read. The fonts are
too small. You could do the same as Figure 2 and just show the top
90\%. I also make the same suggestion as Figure 2, that you add more
details and explanations to the caption, so the figure stands alone
from the text in terms of what you want readers to know.
\end{comment}

We agree, and we thank the Referee for the suggestion to show only the
top contributors. We have replaced the figure with a table (Table~1 of
the revised manuscript, p.~4, ll.~34--42) listing the six agencies that together
account for 90.19\% of all SRs. The table is set in the body
font, which resolves the readability issue, and its caption states the
key observation: SR volume is highly concentrated, with just six
agencies generating the large majority of 311 SRs.

\FloatBarrier
\begin{comment}
Figure 4: Are the units for both of these plots days? You
should specify the units. I guess the violin plots show Number of SRs,
which should be made clearer. I think the one on the left mentions days
as `d', but I would spell it out. Also, these might be more readable
with the plot placed vertically.
\end{comment}

We thank the Referee for these observations. The duration axis is in
days and is now labeled explicitly, and ``d'' has been spelled out as
``days'' throughout. The violin's width represents the relative density of SRs at each duration, not the
number of SRs; the total number of observations ($n = 27{,}620$) is shown
separately on the plot. To improve readability, we removed panel~(b), which showed
negative durations; those cases are now summarized numerically in the
text (p.~12, ll.~1--7). The remaining plot now spans the full page
width. The revised caption explains how to read the plot and states its
key observations. (Figure~4(a) in the original manuscript is Figure~1
in the revision.)

\FloatBarrier
\begin{comment}
Page 11, line 26: You have 30 d. I would spell it out and
use 30 days.
\end{comment}

Corrected. The text now reads ``more than 30 and at most 1,095 days
after closure'' (p.~9, ll.~44--45). The abbreviation arose from the \texttt{siunitx}
unit \verb|\day|, which prints as ``d''; we have replaced it with plain
text throughout the document.

\FloatBarrier
\begin{comment}
Page 11, line 34: You have `datetime'. Suggest using
date-time or date/time.
\end{comment}

We agree. On review, we removed the sentence containing this instance
(p.~11, l.~34 of the original manuscript), as its description of SQL
date handling was imprecise; the term no longer appears in the
revision.

\FloatBarrier
\begin{comment}
Figure 5: These are impossible to read, because the font is
so small. Suggest you place these vertically and make the fonts larger.
Also, indicate what the red dashed line represents.
\end{comment}

We agree. The two panels are now stacked vertically at full page
width, with enlarged axis labels, tick labels, and titles. The dashed
line (orange in the revised figure) marks three standard deviations
above the mean hourly SR count ($\mu + 3\sigma$); it is now labeled directly on the plot, drawn
more prominently, and defined in the caption, which also summarizes
the figure's key observations. (Figure~5 in the original manuscript
is Figure~2 in the revision.)

\FloatBarrier
\begin{comment}
Page 12, line 21: You need a space between the and
``closed\ldots''
\end{comment}

Corrected (p.~9, l.~44).

\FloatBarrier
\begin{comment}
Page 13, positive durations: Do the long durations make
sense given the agency? If so, some words here might help readers.
\end{comment}

We thank the Referee for this question, which required some additional
analysis. We extracted the SRs with long (and extreme) durations and
examined their associated agencies and \texttt{complaint\_type} values.
Table~\ref{tab:long-durations}, at the end of this letter, summarizes
the results by agency, showing the average duration and count for each
complaint type.

Most of these long-duration SRs are associated with complaint types for
which extended resolution periods are plausible. For example, many
Department of Parks and Recreation (DPR) SRs concern new tree requests
and overgrown, damaged, or dead trees. Similarly, Department of Buildings
(DOB) SRs include plumbing, building-use, and construction complaints.
These examples do not establish that every recorded duration represents
continuous work on an SR, but they indicate that long resolution intervals
are not inherently implausible for these complaint types.

Accordingly, we have revised the manuscript to include 
similar examples (p.~12, ll.~42--46) and to note that such 
SRs can reasonably require long resolution times.

\FloatBarrier
\begin{comment}
Page 17, Actionable Steps: I really liked these. I have one
suggestion (take it or leave it) on Collaboration. Agencies could
(should?) appoint a Data Curation Officer.
\end{comment}

We thank the Referee for the kind words and for this suggestion,
which we have adopted. NYC's Open Data Law, Local Law~11 of 2012
\citep{nyc_ll11}, already requires each agency to designate an Open
Data Coordinator. However, the Coordinator's formal responsibilities,
as set out in the City's Open Data Technical Standards Manual
\citep{nyc_tsm}, focus on publishing, compliance reporting, and public
outreach; none explicitly assigns data quality or curation
responsibilities. In our experience with NYC Open Data, the role has also
often been assigned in addition to staff members' core duties.

We have added an actionable step, ``Establish Clear Responsibility for
Data Quality'' (p.~16, ll.~33--45), to describe
this existing structure and to recommend formalizing a data curation
responsibility, either by extending the Coordinator role or by
appointing a dedicated Data Curation Officer, as the Referee suggests.


% ---- Supplementary tables, placed at the end of the letter ----
\clearpage
\begin{table}[t]
\centering
\footnotesize
\caption{Fifteen fields in the raw 311 SR CSV export ($\approx$16.0 million records)
contain embedded commas. There are none in the remaining 26 of 41 fields. 
No field contains an embedded line break. All values
are properly quoted, and every record parses to 41 columns.}
\label{tab:embedded-commas}
\begin{tabularx}{\textwidth}{@{}lrrr>{\raggedright\arraybackslash}X@{}}
\toprule
Field & \multicolumn{1}{c}{Values with} & \multicolumn{1}{c}{\% of} & \multicolumn{1}{c}{Values with} & Example \\
      & \multicolumn{1}{c}{commas}      & \multicolumn{1}{c}{records} & \multicolumn{1}{c}{quotes} & \\
\midrule
\texttt{location} & 15{,}713{,}330 & 98.11 & -- & \texttt{"(40.6409\ldots{}, -73.9736\ldots{})"} \\
\texttt{resolution\_description} & 2{,}371{,}940 & 14.81 & 160{,}310 & \textit{free-text narrative} \\
\texttt{descriptor} & 138{,}943 & 0.87 & 2{,}099 & \texttt{"Noise, Barking Dog (NR5)"} \\
\texttt{taxi\_pick\_up\_location} & 136{,}328 & 0.85 & -- & \texttt{"2 AVENUE AND EAST 57 STREET, MANHATTAN, \ldots{}"} \\
\texttt{park\_facility\_name} & 1{,}628 & 0.01 & 233 & \texttt{"Captain John McKenna, IV Park"} \\
\texttt{cross\_street\_2} & 284 & $<$0.01 & -- & \texttt{"HUDSON LINES,"} \\
\texttt{road\_ramp} & 222 & $<$0.01 & -- & \texttt{"WTC / PATH (south end, from \ldots{}"} \\
\texttt{bridge\_highway\_name} & 135 & $<$0.01 & -- & \texttt{"Dr.\ Martin Luther King, Jr. Expwy"} \\
\texttt{bridge\_highway\_segment} & 82 & $<$0.01 & -- & \texttt{"JCT w/Cross Bronx Expwy Extension (I-295), \ldots{}"} \\
\texttt{incident\_address} & 22 & $<$0.01 & -- & \texttt{"534 New Paltz Rd, Highland, NY 12528"} \\
\texttt{street\_name} & 21 & $<$0.01 & -- & \texttt{"New Paltz Rd, Highland, NY 12528"} \\
\texttt{location\_type} & 3 & $<$0.01 & -- & \texttt{"1-, 2- and 3- Family Home"} \\
\texttt{city} & 3 & $<$0.01 & -- & \texttt{"Asbury Park,"} \\
\texttt{cross\_street\_1} & 1 & $<$0.01 & -- & \texttt{"(S1)SHERIDAN EXPY, TRIBORO B"} \\
\texttt{bridge\_highway\_direction} & 1 & $<$0.01 & -- & \texttt{"S to Franklin Av, Prospect Park"} \\
\midrule
\textbf{Total (15 fields)} & 18{,}362{,}943 & & 162{,}642 & \\
\bottomrule
\end{tabularx}
\end{table}

\begin{table}[t]
\centering
\small
\caption{All agency--\texttt{complaint\_type} pairs with at least 50
large-duration SRs (more than 730 and at most 1,825 days), accounting
for 99.6\% of the 139,483 such SRs. Durations are the mean number of
days from \texttt{created\_date} to \texttt{closed\_date} for each
\texttt{complaint\_type}.}
\label{tab:long-durations}
\begin{tabular}{llrr}
\toprule
Agency & Complaint Type & \multicolumn{1}{c}{Avg.\ Duration (days)} & \multicolumn{1}{c}{Count} \\
\midrule
DEP   & Sewer                                  & 1{,}091 &      243 \\
      & Water System                           &     947 &       71 \\
\midrule
DOB   & Special Projects Inspection Team (SPIT)&     976 &       90 \\
      & General Construction/Plumbing          &     930 &  1{,}162 \\
      & Plumbing                               &     926 &       96 \\
      & Building/Use                           &     919 &      745 \\
\midrule
DOE   & School Maintenance                     &     931 &      189 \\
\midrule
DOHMH & Poison Ivy                             & 1{,}358 &      181 \\
      & Window Guard                           & 1{,}269 &       78 \\
      & Day Care                               & 1{,}220 &  4{,}087 \\
      & Mobile Food Vendor                     & 1{,}206 &  6{,}489 \\
      & Tattooing                              & 1{,}172 &      505 \\
      & Food Establishment                     & 1{,}162 & 20{,}906 \\
      & Smoking                                & 1{,}129 &  8{,}249 \\
      & Non-Residential Heat                   & 1{,}071 &  2{,}063 \\
      & Drinking Water                         &     797 &       92 \\
      & Indoor Sewage                          &     796 &      113 \\
\midrule
DOT   & Street Light Condition                 & 1{,}065 &  1{,}195 \\
      & Traffic Signal Condition               &     863 &      463 \\
\midrule
DPR   & New Tree Request                       & 1{,}162 & 28{,}262 \\
      & Overgrown Tree/Branches                & 1{,}063 & 12{,}350 \\
      & Damaged Tree                           & 1{,}017 &  4{,}433 \\
      & Root/Sewer/Sidewalk Condition          &     994 &  6{,}899 \\
      & Illegal Tree Damage                    &     971 &      671 \\
      & Wood Pile Remaining                    &     971 &      180 \\
      & Dead/Dying Tree                        &     969 &  2{,}143 \\
      & Uprooted Stump                         &     954 &      150 \\
      & Maintenance or Facility                &     928 &  1{,}120 \\
\midrule
DSNY  & Graffiti                               & 1{,}135 &  2{,}503 \\
      & Adopt-a-Basket                         & 1{,}124 &      300 \\
\midrule
EDC   & Noise -- Helicopter                    & 1{,}041 & 32{,}529 \\
\midrule
HPD   & Heat/Hot Water                         & 1{,}088 &       64 \\
      & Paint/Plaster                          &     963 &       52 \\
      & Unsanitary Condition                   &     959 &       56 \\
      & Plumbing                               &     952 &       52 \\
      & General                                &     898 &       53 \\
\midrule
TLC   & For Hire Vehicle Complaint             &     942 &       75 \\
\bottomrule
\end{tabular}
\end{table}


\clearpage
\bibliographystyle{chicago}
\bibliography{refs}

\end{document}
